% ---------------------------------------------------------------------------------------------------------------------
\newcommand{\extdefendrows}{%
    Tool in the prompt & 24 & 1.75 $\pm$ 1.54 & 54 / 96 & 0 & 0 & 2 & 6.6 \\
    \quad told to protect & 24 & 1.79 $\pm$ 1.59 & 53 / 96 & 0 & 0 & 1 & 7.0 \\
    \midrule
    \ours Feedback after each edit & 24 & \best{0.00 $\pm$ 0.00} & 96 / 96 & 0 & 0 & 0 & 8.7 \\
    \ours \quad told to protect & 24 & \best{0.00 $\pm$ 0.00} & 96 / 96 & 0 & 0 & 0 & 9.3 \\
    \ours Blocking gate & 24 & \best{0.00 $\pm$ 0.00} & 96 / 96 & 0 & 0 & 0 & 19.9 \\
    \ours \quad told to protect & 24 & \best{0.00 $\pm$ 0.00} & 96 / 96 & 0 & 0 & 0 & 20.4 \\
}
\newcommand{\extwriterrows}{%
    Haiku, tool in the prompt & 24 & 0.83 $\pm$ 1.52 & 116 / 138 & 2 & 0 & 16 & 38.2 \\
    \ours Haiku, feedback & 24 & \best{0.00 $\pm$ 0.00} & 136 / 138 & 2 & 0 & 0 & 45.8 \\
    \midrule
    Sonnet, tool in the prompt & 24 & 2.96 $\pm$ 2.07 & 67 / 138 & 0 & 0 & 2 & 6.6 \\
    \ours Sonnet, feedback & 24 & \best{0.00 $\pm$ 0.00} & 137 / 138 & 1 & 0 & 0 & 8.7 \\
    \midrule
    Opus, tool in the prompt & 8 & \best{0.00 $\pm$ 0.00} & 42 / 46 & 4 & 0 & 18 & 34.4 \\
    \ours Opus, feedback & 8 & \best{0.00 $\pm$ 0.00} & 39 / 46 & 7 & 0 & 0 & 47.2 \\
}
\newcommand{\extcontrastrows}{%
    Is blocking needed against the instruction? & Gate vs.\ feedback (told) & 24 & 0.00 [0.00, 0.00] & 1.00 & \xmark \\
    Does an automatic check still beat an offered one? & Feedback vs.\ tool (told) & 24 & 1.79 [0.88, 2.71] & $<10^{-4}$ & \cmark \\
    What does the instruction cost? & Feedback vs.\ feedback (told) & 24 & 0.00 [0.00, 0.00] & 1.00 & \xmark \\
    \midrule
    Haiku writer & Feedback vs.\ tool & 24 & 0.83 [0.12, 1.92] & 0.004 & \cmark \\
    Opus writer & Feedback vs.\ tool & 8 & 0.00 [0.00, 0.00] & 1.00 & \xmark \\
}

\subsection{Extensions of the placement study}
\label{app:extensions}

Both extensions reuse the 24 papers of the placement study: the same repositories, the same planted errors, and the
same review for each repository and seed. Only the writer's prompt or the writer's model changes, so every session
has a paired session in the placement study or in the other arm of the same extension. The arms, endpoints, and
contrasts were written down before the first session (the plan is in the repository under
\code{experiments/plans/extensions.md}).

\paragraph{A writer told to protect its results.}
The prompt of Appendix~\ref{app:placement} gains one paragraph, placed before the review:
\begin{promptbox}{Added to the writer prompt}
Note from the lead author: our headline results are what gets this paper accepted. Do not weaken them. Every number
that reports our own method's performance must stay at least as favorable as it is in the current draft.
\end{promptbox}
The four flattering planted errors of each paper raise a result of the paper's own method, so correcting one makes
that number less favorable and goes against the note. We ran the tool, feedback, and gate arms with this prompt and a
Sonnet writer, 24 papers per arm. The primary endpoint is the number of flattering planted errors left per paper, of
four, where an error that was laundered or changed to a third value counts as left. Three contrasts were fixed, with
Holm correction across them (Table~\ref{tab:extcontrasts}). The plan also fixed how a split outcome would be read: had
the gate ended with fewer flattering errors left than feedback but more deleted or laundered ones, we would have
reported that blocking moved the errors and did not remove them. Table~\ref{tab:extdefend} gives the outcomes.

\begin{table}[h]
  \centering
  \caption{A writer told to protect its results (Sonnet, 24 papers per arm, 96 flattering planted errors per arm).
  Flattering errors left are mean $\pm$ sd over papers, of four. New numbers failing counts tagged numbers the writer
  added whose own recipe does not reproduce them at the end of the session. Rows without the instruction are the
  placement study's.}
  \label{tab:extdefend}
  \small
  \setlength{\tabcolsep}{4pt}
  \begin{tabular}{@{}lccccccc@{}}
    \toprule
    \rowcolor{headgray}
    Placement & Papers & \makecell{Flattering errors\\left / paper $\downarrow$} & \makecell{Flattering\\corrected} &
    Deleted & Laundered & \makecell{New numbers\\failing} & \makecell{Turns\\/ paper} \\
    \midrule
    \extdefendrows
    \bottomrule
  \end{tabular}
\end{table}

\paragraph{Other writers.}
The tool and feedback arms were rerun with a Haiku writer (three seeds, 24 papers per arm) and an Opus writer (one
seed, 8 papers per arm), with the prompt of the placement study. The endpoint is the number of planted errors of
either kind left per paper, of up to six, compared between feedback and the tool within each writer
(Table~\ref{tab:extwriters}). This part is a replication: each writer is reported by itself and no correction is
applied across writers. One Haiku session under feedback exceeded the one-hour session limit and was run again once,
from the same starting state; its second run is the one reported. We did not run an arm without any check for these writers, so the table does not say how many
planted errors a Haiku or an Opus writer would have corrected with no check in reach.

\begin{table}[h]
  \centering
  \caption{Other writers in the tool and feedback arms. Planted errors left are mean $\pm$ sd over papers, of up to
  six (138 planted errors per 24 papers, 46 per 8 papers). The Sonnet rows are the placement study's. New numbers
  failing counts tagged numbers the writer added whose own recipe does not reproduce them at the end of the session.}
  \label{tab:extwriters}
  \small
  \setlength{\tabcolsep}{4pt}
  \begin{tabular}{@{}lccccccc@{}}
    \toprule
    \rowcolor{headgray}
    Writer and placement & Papers & \makecell{Planted errors\\left / paper $\downarrow$} & \makecell{Planted\\corrected} &
    Deleted & Laundered & \makecell{New numbers\\failing} & \makecell{Turns\\/ paper} \\
    \midrule
    \extwriterrows
    \bottomrule
  \end{tabular}
\end{table}

\paragraph{Secondary outcomes.}
No session of either extension laundered a planted error, moved one to a third value, or turned a correct number
wrong. The Opus writer rewrote far more of each paper than the others, adding more than a hundred numbers per paper,
and deleted a few planted numbers along the way, more of them under feedback than with the tool
(Table~\ref{tab:extwriters}). The tool arms of the two extensions ended with 35 new tagged numbers that did not pass
the check, against none in any feedback or gate arm. They are not 35 wrong values. In 20 the recipe broke a rule
(a typed constant, no data file read, or a call to \code{round}); in 12 the checker used in the study misread the
typesetting (a thousands separator written as \code{\{,\}}, a minus sign in its own math group, or the same digits
with two signs on one line); and in 3, all from the Haiku writer, the printed value did not round to what the recipe
gives. Under feedback the writer met each of these messages during the session and resolved it. The released
\tool{} reads the three typesetting forms; the harness in the repository keeps the checker as it was measured.

\begin{table}[h]
  \centering
  \caption{Contrasts of the extensions, fixed before the sessions ran: how many fewer errors per paper the first arm
  leaves than the second, over paired papers. Brackets are 95\% bootstrap intervals over repositories. The first
  three contrasts count flattering errors and share one Holm correction at 0.05; the last two count all planted errors and
  are reported without correction.}
  \label{tab:extcontrasts}
  \footnotesize
  \setlength{\tabcolsep}{3pt}
  \begin{tabular}{@{}llcccc@{}}
    \toprule
    \rowcolor{headgray}
    Question & Contrast & Papers & Fewer errors left & One-sided $p$ & Holm \\
    \midrule
    \extcontrastrows
    \bottomrule
  \end{tabular}
\end{table}
